\section*{Capítulo \ref{ch:b3:quantum-model-systems} --- Mecânica quântica para químicos: sistemas-modelo}

\begin{solution}{exo:b3:quantum-model-systems:1}
$E_1 = h^2/8m_eL^2$. (a) $L = \qty{1.00}{nm}$: $E_1 = \qty{6.02e-20}{J} =
\qty{0.376}{eV}$; $\Delta E_{12} = 3E_1$ e $\lambda = hc/3E_1 =
\qty{1099}{nm}$, no infravermelho próximo. (b) $L = \qty{0.50}{nm}$: $E_1$ é quatro
vezes maior, \qty{2.41e-19}{J} $= \qty{1.50}{eV}$, e $\lambda =
\qty{275}{nm}$, no ultravioleta.
\end{solution}

\begin{solution}{exo:b3:quantum-model-systems:2}
$n_x^2 + n_y^2 + n_z^2 = 3$ (1,1,1): $g = 1$; $6$ (2,1,1 e permutações):
$g = 3$; $9$ (2,2,1): $g = 3$; $11$ (3,1,1): $g = 3$; $12$ (2,2,2): $g = 1$.
\end{solution}

\begin{solution}{exo:b3:quantum-model-systems:3}
$[\hat x^2,\hat p]f = -\iu\hbar x^2f' + \iu\hbar(x^2f)' = 2\iu\hbar xf$, logo
$[\hat x^2,\hat p] = 2\iu\hbar\hat x$. $[\hat p^2,\hat x]f = -\hbar^2[(xf)'' -
xf''] = -2\hbar^2f'$, logo $[\hat p^2,\hat x] = -2\iu\hbar\hat p$.
\end{solution}

\begin{solution}{exo:b3:quantum-model-systems:4}
$E_J/hc = BJ(J+1)$: 0, 21.19, 63.56 e \qty{127.12}{cm^{-1}}, com $g = 1$,
3, 5, 7. $I = h/(8\pi^2cB) = \qty{2.643e-47}{kg.m^2}$ (com $c$ em
\unit{cm/s}).
\end{solution}

\begin{solution}{exo:b3:quantum-model-systems:5}
$|\psi_n|^2$ é simétrica em relação a $L/2$, logo $\langle x\rangle = L/2$. Com
$\sin^2u = \frac12(1 - \cos2u)$ e duas integrações por partes,
$\langle x^2\rangle = L^2\big(\frac13 - \frac{1}{2n^2\pi^2}\big)$. Para $n = 1$,
$\langle x^2\rangle = 0.283L^2$, menos que o valor clássico $L^2/3$: o estado
fundamental se acumula no meio. À medida que $n$ cresce, recupera-se o valor clássico.
\end{solution}

\begin{solution}{exo:b3:quantum-model-systems:6}
$1\,\unit{cm^{-1}} = \qty{0.011963}{kJ/mol}$. EPZ(\ce{H2}) $= 2200.6 \times
0.011963 = \qty{26.33}{kJ/mol}$, EPZ(\ce{D2}) $= \qty{18.63}{kJ/mol}$; diferença
\qty{7.69}{kJ/mol}. Mesma \omterm{def:b3:quantum-model-systems:oscillator}{constante de força}, logo $\tilde\omega \propto \mu^{-1/2}$
e a razão esperada é $\sqrt{\mu_D/\mu_H} = 1.4137$ (massas atômicas
2.01410 e 1.00783); a razão medida é 1.4127, com diferença de 0.07\,\% (a
pequena diferença vem dos elétrons, que não acompanham os núcleos
perfeitamente).
\end{solution}

\begin{solution}{exo:b3:quantum-model-systems:7}
$\int_0^Lx^2(L - x)^2\,\dd x = L^5/30$, logo $N = \sqrt{30/L^5}$. Então
\[
\langle\psi_1|\phi\rangle = \sqrt{\tfrac2L}\sqrt{\tfrac{30}{L^5}}\int_0^Lx(L -
x)\sin\tfrac{\pi x}{L}\,\dd x = \frac{\sqrt{60}}{L^3}\cdot\frac{4L^3}{\pi^3} =
\frac{4\sqrt{60}}{\pi^3} = 0.99928 .
\]
A parábola é 99.86\,\% estado fundamental (o quadrado da sobreposição).
\end{solution}

\begin{solution}{exo:b3:quantum-model-systems:8}
$\langle f|\hat pg\rangle = -\iu\hbar\int f^*g'\,\dd x = -\iu\hbar[f^*g] +
\iu\hbar\int f^{*\prime}g\,\dd x = \int(-\iu\hbar f')^*g\,\dd x = \langle\hat
pf|g\rangle$, já que o colchete se anula nas extremidades. Para $\hat x\hat p$, usando que $\hat x$ e $\hat p$ são, cada um, \omterm{def:b3:quantum-model-systems:hermitian}{hermitianos}:
$\langle f|\hat x\hat pg\rangle = \langle\hat xf|\hat pg\rangle = \langle\hat p\hat
xf|g\rangle$, e $\hat p\hat x = \hat x\hat p - \iu\hbar \ne \hat x\hat p$: não é
\omterm{def:b3:quantum-model-systems:hermitian}{hermitiano} (sua forma simetrizada $\frac12(\hat x\hat p + \hat p\hat x)$ é).
\end{solution}

\begin{solution}{exo:b3:quantum-model-systems:9}
$L = 6 \times 139.7 = \qty{838}{pm}$, $N = 6$: $\lambda = 8m_ecL^2/(7h) =
\qty{331}{nm}$. A absorção está no ultravioleta: o hexatrieno é
incolor.
\end{solution}

\begin{solution}{exo:b3:quantum-model-systems:10}
$\kappa = \sqrt{2m(V_0 - E)}/\hbar$ com $V_0 - E = \qty{0.20}{eV}$:
$\kappa_H = \qty{9.82e10}{m^{-1}}$, $\kappa_D = \qty{1.388e11}{m^{-1}}$. Os
prefatores se cancelam na razão: $T_H/T_D = \eu^{2a(\kappa_D - \kappa_H)} =
\eu^{4.06} = 58$, o valor exato: a barreira é espessa para os dois
($\kappa a = 4.9$ e 6.9).
\end{solution}

\begin{solution}{exo:b3:quantum-model-systems:11}
Níveis $E_m = m^2\hbar^2/2m_eR^2$: $m = 0$ comporta dois elétrons, $m = \pm1$
quatro. O nível preenchido mais alto é $|m| = 1$, o vazio mais baixo $|m| = 2$: $\Delta E =
3\hbar^2/2m_eR^2 = \qty{9.38e-19}{J}$, $\lambda = \qty{212}{nm}$, no
ultravioleta, como a forte absorção do benzeno.
\end{solution}

\begin{solution}{exo:b3:quantum-model-systems:12}
$\langle r\rangle = \frac{4\pi}{\pi a^3}\int_0^\infty r^3\eu^{-2r/a}\dd r =
\frac{4}{a^3}\cdot\frac{3!}{(2/a)^4} = \frac32a$, \qty{79.4}{pm} para $a = a_0$.
$\langle 1/r\rangle = \frac{4}{a^3}\cdot\frac{1!}{(2/a)^2} = \frac1a$, logo
$\langle V\rangle = -e^2/(4\pi\varepsilon_0a)$. Como $E_1 =
-e^2/(8\pi\varepsilon_0a)$ (de $a = 4\pi\varepsilon_0\hbar^2/\mu e^2$),
$\langle V\rangle = 2E_1$, e a energia cinética média é $-E_1$.
\end{solution}

\begin{solution}{pb:b3:quantum-model-systems:1}
\textbf{1.} Cada um dos $j$ átomos fornece um orbital p; os $j - 1$ carbonos e
um nitrogênio fornecem um elétron cada, e o outro nitrogênio seu par não
ligante, pois a carga do cátion se distribui pela cadeia: $N = j + 1$, ou seja, 6,
8 e 10.
\textbf{2.} Seis elétrons preenchem $n = 1$, 2, 3: HOMO $n = 3$, LUMO $n = 4$.
\textbf{3.} $\Delta E = E_{N/2+1} - E_{N/2} = (N+1)h^2/8mL^2$ e $\lambda =
hc/\Delta E = 8mcL^2/h(N+1)$.
\textbf{4.} $L_0 = 4l$, $6l$, $8l$: 559, 838 e \qty{1118}{pm}.
\textbf{5.} 147, 257 e \qty{374}{nm}.
\textbf{6.} Todos curtos demais (por fatores 3.6, 2.3 e 1.9): a caixa do modelo é
pequena demais, já que $\lambda \propto L^2$; os elétrons se movem por uma
distância maior que a cadeia entre os nitrogênios.
\textbf{7.} Com a origem no centro, $\psi_n \propto \cos(n\pi u/L)$ para
$n$ ímpar e $\sin(n\pi u/L)$ para $n$ par ($u = x - L/2$): funções pares e
ímpares de $u$.
\textbf{8.} Se $n + n'$ é par, $\psi_n$ e $\psi_{n'}$ têm a mesma paridade;
seu produto multiplicado por $u$ é ímpar e sua integral sobre um intervalo
simétrico é nula.
\textbf{9.} O HOMO $N/2$ e o LUMO $N/2 + 1$ têm $n + n' = N + 1$, ímpar: sempre
permitida.
\textbf{10.} Segundo corante: HOMO 4, LUMO 5. $4 \to 6$: soma par, proibida.
$3 \to 5$: soma par, proibida.
\textbf{11.} A próxima transição permitida a partir de $n = 4$ é $4 \to 7$, com
$\Delta E$ maior por um fator $(49 - 16)/(25 - 16) = 3.67$: em $\lambda/3.67$,
\qty{164}{nm} para o segundo corante, no ultravioleta profundo.
\textbf{12.} Ela depende apenas da simetria do potencial em relação ao
centro da cadeia, que os corantes simétricos reais compartilham.
\textbf{13.} $E = hc/\lambda = \qty{3.294e-19}{J} = \qty{2.06}{eV}$.
\textbf{14.} $hc\tilde\omega_e = \qty{0.371}{eV}$.
\textbf{15.} $2hcB = \qty{4.21e-22}{J} = \qty{2.63}{meV}$.
\textbf{16.} $kT = \qty{25.7}{meV}$.
\textbf{17.} Apenas os níveis rotacionais ($2hcB \ll kT$); as excitações
vibracionais ($15kT$) e eletrônicas ($80kT$) praticamente não estão povoadas.
\textbf{18.} Eletrônicas: visível e ultravioleta; vibracionais: infravermelho;
rotacionais: micro-ondas (e infravermelho distante).
\textbf{19.} $L = \sqrt{\lambda h(N+1)/8mc}$ com $N = 8$: $L = \qty{1283}{pm}$.
\textbf{20.} $\delta = (1283 - 838)/2 = \qty{222}{pm}$.
\textbf{21.} $L = 559 + 445 = \qty{1003}{pm}$, $\lambda = \qty{474}{nm}$, 9.5\,\%
abaixo de 524 nm.
\textbf{22.} $L = 1118 + 445 = \qty{1562}{pm}$, $\lambda = \qty{732}{nm}$, 2.9\,\%
acima de 711 nm.
\textbf{23.} Para dentro dos anéis aromáticos que carregam os nitrogênios: o
sistema $\pi$ não termina nos átomos de nitrogênio, e a caixa efetiva inclui
parte de cada anel.
\textbf{24.} $\delta/l = 222/139.7 = 1.6$: \textbf{a caixa se estende cerca de
\qty{222}{pm}, uma ligação e meia, além de cada nitrogênio}, e com esse único
comprimento o modelo reproduz as três cores com erro inferior a 10\,\%.
\end{solution}
